\subsection{Random-endpoint vertex cover: an optional expectation proof}
\label{app:randomvc}
For an unweighted graph, repeatedly choose any uncovered edge and select one of its two
endpoints uniformly at random, then remove all edges incident to the selected vertex.
Return the selected vertices. Each iteration removes at least its chosen edge, and every
removed edge is covered, so the output is feasible.

Fix a minimum cover $C^*$. There are at most $n$ iterations, since selected vertices are
distinct. For each $i\in\{1,\ldots,n\}$, let $A_i$ indicate that iteration $i$ occurs,
and let $Y_i$ indicate that it selects a vertex of $C^*$; after termination set both to zero.
Conditional on the history of any active iteration, the chosen edge has at least one
endpoint in $C^*$, so
\[
\mathbb E[Y_i\mid\text{history}]\ge\tfrac12 A_i.
\]
Taking expectations and summing gives
\[
\tfrac12\mathbb E[|C|]=\tfrac12\sum_i\mathbb E[A_i]
\le\mathbb E\!\left[\sum_iY_i\right]\le |C^*|.
\]
Thus $\mathbb E[|C|]\le2\OPT$. This does not claim a factor-2 bound for every run.
Maintaining active edges and removing a vertex's incident edges gives polynomial time.
The simpler maximal-matching algorithm in the main notes is deterministic and has the
same approximation factor.

\section{Source Notes and Scope}
\label{app:sources}
The core material follows the supplied L1--L4 approximation/LP decks, Network Flow I--III,
and Tutorials 1--4. The merged tutorial PDF includes Tutorial 4 on PDF pages 21--28.
Network Flow III slide 36 announces the Week 6 cutoff. The sample finals are used for
practice style; their exam instructions do not determine the midterm's rules.

\begin{itemize}
\item \textbf{L1, slide 23:} knapsack's capacity constraint is $\sum_i w_ix_i\le W$.
The opposite sign printed there is inconsistent with the problem and surrounding slides.
\item \textbf{L4-LP, slide 21:} the displayed clause expression is a lower bound in
general, and the final comparison with $(1-1/e)y_i^*$ is $\ge$, not $\le$.
\item \textbf{Network Flow III, scaling pseudocode:} initialize the threshold to a power
of two, or otherwise ensure threshold 1 is processed. Real halving from an arbitrary
integer maximum capacity can skip it. Describe flow algorithms by their path-selection
rules as well as their names, since labels differ across sources.
\item \textbf{AY25/26 final, solution Q4C:} the assignment arc in its layered network
is $(u_{i,j},v)$, not $(u_i,v)$. If a schedule covers a day redundantly and all employee
constraints are upper bounds, discard excess assignments before constructing a unit
flow per day.
\item \textbf{AY25/26 final, solution Q5A:} checking only conflicting pairs and triangles
does not suffice for same-side/opposite-side constraints. A 5-cycle of opposite-side
constraints is infeasible but has no triangle. Correct feasibility propagates parity by
BFS/DFS across every constraint edge, rejecting any inconsistent assignment. The SDP
parts of that question are not developed in the supplied Week 1--6 decks.
\item \textbf{AY24/25 final, solution Q6:} strong LP duality equates the fractional
matching/cover optima. Identifying these values with the integral combinatorial problems
also uses the integrality property for bipartite graphs. Matching via flow is part of the
supplied scope; the LP-duality argument is additional material.
\end{itemize}

